\documentclass[12pt]{article}
\usepackage[utf8]{inputenc}
\usepackage[T1]{fontenc}
\usepackage{hyperref}
\title{Yawiyo Language Resources for K-12 Curriculum Development}
\author{Compiled by Claude Sonnet 5.5 (claude-sonnet-5-5) for Andr\'as Kornai}
\date{}
\begin{document}
\maketitle

\section{Introduction}
Yawiyo (also called Yabio; ISO 639-3: ybx) is spoken in East Sepik Province, Papua New Guinea. Wikipedia cites 66 speakers (2000 census). Joshua Project places the language area along the upper Wano river, spanning East Sepik and Sandaun Provinces. The Endangered Languages Project rates the language severely endangered and cites Conrad and Lewis (1988) for a population of 60 to 100. Sources disagree on classification: Wikipedia says Sepik, while the Endangered Languages Project lists it as unclassified.

\section{Primary Educational Resources}
The search found no primers, textbooks or school programmes for Yawiyo.

\section{Linguistic and Reference Materials}
\begin{itemize}
\item Wikipedia overview: \url{https://en.wikipedia.org/wiki/Yawiyo_language}
\item Endangered Languages Project entry: \url{https://www.endangeredlanguages.com/lang/3212?hl=fr}
\end{itemize}
The search found no grammar or dictionary of Yawiyo.

\section{Religious and Community Resources}
\begin{itemize}
\item Joshua Project profile (listed as Yabio): \url{https://joshuaproject.net/people_groups/15932}
\end{itemize}

\section{Digital and Audio Resources}
The search found no audio or digital resources specific to Yawiyo.

\section{Curriculum Development Context}
With about 66 speakers and no known teaching materials, any curriculum work would begin with documentation. That would mean recording elders, writing down an orthography and building a small primer and word list. Tok Pisin and English would remain the main school languages.

\end{document}